\subsection{局部闭子空间}

定义 2. 称拓扑空间 X 的子集 L 在点 \(x \in L\) 处是局部闭的，如果存在 x 在 X 中的一个邻域 V，使得 \(L \cap V\) 是子空间 V 的一个闭子集。若 L 在每个 \(x \in L\) 处都是局部闭的，则称 L 在 X 中是局部闭的。

注. 设 F 是 X 的一个子集，使得对 X 的每个点 x 存在 x 的一个邻域 V 满足 \(V \cap F\) 在子空间 V 中是闭集；则由第 1 目命题 3，F 在 X 中是闭集。另一方面，下面的命题 5 表明，一般存在在 X 中不闭的局部闭集。

命题 5. 对拓扑空间 X 的子集 L 而言，下列诸性质等价：

\begin{enumerate}
\def\labelenumi{\alph{enumi})}
\item
  L 在 X 中是局部闭的。
\item
  L 是 X 的一个开子集与一个闭子集的交。
\item
  L 在它在 X 中的闭包 \(\overline{L}\) 中是开集。
\end{enumerate}

若 L 局部闭，则对每个 \(x \in L\)，存在 x 在 X 中的一个开邻域 \(V_{x}\)，使得 \(L \cap V_{x}\) 在 \(V_{x}\) 中是闭集；\(U = \bigcup_{x \in L} V_{x}\) 在 X 中是开集，且第 1 目命题 3 表明 L 在 U 中是闭集；因此 a) 蕴含 b)。若 \(L = U \cap F\)，其中 U 是开集而 F 在 X 中是闭集，则有 \(\overline{L} \subset F\)；于是 \(L \subset U \cap \overline{L} \subset U \cap F = L\)，这表明 \(L = U \cap \overline{L}\) 在 \(\overline{L}\) 中是开集，故 b) 蕴含 c)。最后，若 \(L = U \cap \overline{L}\)，其中 U 在 X 中是开集，则 \(\overline{L}\) 在 U 中是闭集，从而是局部闭的，于是 c) 蕴含 a)。

推论. 设 \(f: \mathbf{X} \to \mathbf{X}'\) 是一个连续映射，\(\mathbf{L}'\) 是 \(\mathbf{X}'\) 的一个局部闭子集；则 \(\overline{f}^{1}(\mathbf{L}')\) 在 \(\mathbf{X}\) 中是局部闭的。

这由上面的命题 5 以及 § 2 第 1 目的定理 1 立即得出。

